\chapter{Problems to Work On}
\label{ch:problems}

\begin{center}\itshape
A good problem is one you can start on today\\
and still not have finished in ten years.\\
Here are seven, with the first steps already taken.
\end{center}

Every problem in this chapter is open. Every one of them can be stated to a curious teenager, has a history with real people in it, and has a partial result that we prove in full. The partial
results are not the hard part of the problem; they are the place where a beginner stands before taking a step of their own. The chapter ends with advice on choosing a problem, and on where to find more.

A warning, given in the same spirit as the problems: none of these is a good choice for a first research project if the goal is to solve it. They are a good choice if the goal is to learn what
an open problem feels like from the inside, to find the special cases nobody has written down, and to compute. Chapter~\ref{ch:projects} gives projects with more modest and more reachable aims.

\section{Singmaster's conjecture}

Pascal's triangle is full of repetitions. The number $10$ appears four times: $\binom{10}1=\binom{10}9=\binom53=\binom52$. The number $3003$ appears eight times. In 1971 David Singmaster asked how often a number can
appear. Let $N(a)$ be the number of pairs $(n,k)$ with $0\le k\le n$ and $\binom nk=a$, for an integer $a\ge2$.

\begin{proposition}[Singmaster, 1971]\label{prop:singmaster}
$N(a)\le2\log_2a+2$.
\end{proposition}

\begin{proof}
By symmetry it suffices to count the pairs with $k\le n/2$ and multiply by $2$. If $1\le k\le n/2$, then $\binom nk\ge\binom{2k}k=\prod_{i=1}^k\frac{k+i}i\ge2^k$ (for fixed $k$, $\binom nk$ increases with $n$, and each factor $\frac{k+i}i$ is at least $2$). So $\binom nk=a$ forces $k\le\log_2a$; and $k=0$ gives $1\ne a$. For each fixed $k\ge1$, $n\mapsto\binom nk$ is strictly increasing on $n\ge k$, so at most one $n$ works. There are at most $\log_2a+1$ values of $k$, so at most that many pairs with $k\le n/2$, and $N(a)\le2\log_2a+2$.
\end{proof}

Singmaster conjectured that $N(a)$ is in fact bounded. The best known upper bound, due to Daniel Kane (2007), is
\[
N(a)=O\Bigl(\frac{\log a\,\log\log\log a}{(\log\log a)^3}\Bigr),
\]
and the conjecture is open. In the other direction, Singmaster found in 1975 that infinitely many numbers appear at least six times. Write $F_m$ for the Fibonacci numbers ($F_1=F_2=1$).

\begin{proposition}[Singmaster, 1975]\label{prop:singmasterfamily}
For $i\ge1$ let $n=F_{2i+2}F_{2i+3}-1$ and $k=F_{2i}F_{2i+3}-1$. Then $\binom{n+1}{k+1}=\binom n{k+2}$, and this number appears at least six times in Pascal's triangle.
\end{proposition}

\begin{proof}
Since $\binom{n+1}{k+1}=\frac{n+1}{k+1}\binom nk$ and $\binom n{k+2}=\frac{(n-k)(n-k-1)}{(k+1)(k+2)}\binom nk$, the equality is equivalent to $(n+1)(k+2)=(n-k)(n-k-1)$. Write $a=F_{2i}$, $b=F_{2i+1}$, $c=F_{2i+2}$, $e=F_{2i+3}$. Then $n+1=ce$, $k+2=ae+1$, and $n-k=(c-a)e=be$. The equation becomes $ce(ae+1)=be(be-1)$, that is $c(ae+1)=b(be-1)$, that is $e(b^2-ac)=b+c=e$. So it is equivalent to $b^2-ac=1$, which is Cassini's identity $F_{2i}F_{2i+2}-F_{2i+1}^2=-1$. (Cassini's identity $F_{m-1}F_{m+1}-F_m^2=(-1)^m$ is the determinant of the $m$-th power of $\left(\begin{smallmatrix}1&1\\1&0\end{smallmatrix}\right)$, whose entries are $F_{m+1},F_m,F_m,F_{m-1}$.)

The common value $a'$ appears at $(n+1,k+1)$, $(n+1,n-k)$, $(n,k+2)$, $(n,n-k-2)$, and at $(a',1)$, $(a',a'-1)$. These six positions are distinct. Indeed $k+1\ge4$, and $k+2<\frac n2$, because $n-2(k+2)=(F_{2i+2}-2F_{2i})F_{2i+3}-3=F_{2i-1}F_{2i+3}-3\ge5-3$; so the four positions in rows $n$ and $n+1$ are distinct and not at the edges, and $a'>n+1$ puts the last two in another row.
\end{proof}

For $i=1$: $n=14$, $k=4$, and $\binom{15}5=\binom{14}6=3003$, which also equals $\binom{78}2$, so it appears eight times. No other number with eight appearances is known, and no number with more.

\section{The lonely runner}

Runners start together from the same point of a circular track of length $1$ and run forever at constant, distinct speeds. A runner is \emph{lonely} at a moment when every other runner is at distance at least $\frac1n$ along the track, where $n$ is the number of runners. The conjecture, formulated by J\"org Wills in 1967 and independently by Thomas Cusick in 1973 (in the language of view obstruction), and named by Luis Goddyn, says that each runner is lonely at some moment. Making one runner stand still (subtract its speed from all), it becomes: \emph{for any $k$ distinct positive integers $v_1,\dots,v_k$, there is a real $t$ with $\|v_it\|\ge\frac1{k+1}$ for all $i$}, where $\|x\|$ is the distance from $x$ to the nearest integer. (Real speeds reduce to integer ones; we take that reduction for granted.)

\begin{proposition}\label{prop:lonely2}
The conjecture holds for $k=1$ and $k=2$.
\end{proposition}

\begin{proof}
For $k=1$, take $t=\frac1{2v_1}$: then $\|v_1t\|=\frac12$. For $k=2$, let $v_1=a$, $v_2=b$, and divide both by $\gcd(a,b)$, which only rescales time; so $\gcd(a,b)=1$, and $a+b\ge3$. For $t=\frac x{a+b}$ with $x$ an integer, $at+bt=x$, so $\|at\|=\|bt\|$. Since $\gcd(a,a+b)=\gcd(a,b)=1$, as $x$ runs over $0,\dots,a+b-1$, $ax$ runs over all residues modulo $a+b$. The interval $[\frac{a+b}3,\frac{2(a+b)}3]$ has length at least $1$, so it contains an integer $r$; choose $x$ with $ax\equiv r\pmod{a+b}$. Then $\|at\|=\|\frac r{a+b}\|\ge\frac13$, and $\|bt\|$ is the same.
\end{proof}

The case $k=3$ was proved by Betke and Wills (1972) and by Cusick (1974); the case of seven runners by Barajas and Serra in 2008, after work of Bohman, Holzman, Kleitman and Renault on six. The bound $\frac1{k+1}$ cannot be improved: for speeds $1,2,\dots,k$ it is attained (Exercise~\ref{ex:lonelysharp}).

\section{Graceful trees}

A tree with $n$ edges has $n+1$ vertices. A \emph{graceful labelling} gives the vertices distinct labels in $\{0,1,\dots,n\}$ such that the $n$ differences $|\ell(u)-\ell(v)|$ along the edges are exactly $1,2,\dots,n$. The \emph{graceful tree conjecture}, from the work of Gerhard Ringel and Anton Kotzig in the 1960s and formalised by Alexander Rosa in 1967, says that every tree has one. It is the most famous of the graph labelling problems, and Rosa's motivation was a decomposition problem: a graceful labelling of a tree $T$ with $n$ edges yields a decomposition of the complete graph $K_{2n+1}$ into $2n+1$ copies of $T$ (Exercise~\ref{ex:ringel}).

A \emph{caterpillar} is a tree in which removing the leaves leaves a path (the \emph{spine}). Paths and stars are caterpillars.

\begin{lemma}\label{lem:gracefulorder}
Let $T$ be a tree with $n$ edges whose vertices are coloured top and bottom so that every edge joins a top and a bottom vertex. Suppose the edges can be ordered $e_1,\dots,e_n$ so that for each $m\ge2$, $e_m$ shares an endpoint with $e_{m-1}$ and its other endpoint is new (not in $e_1,\dots,e_{m-1}$). Then $T$ is graceful.
\end{lemma}

\begin{proof}
Number the top vertices $0,1,2,\dots$ and the bottom vertices $0,1,2,\dots$ in the order in which they first appear in $e_1,e_2,\dots$; label the top vertex number $i$ by $i$ and the bottom vertex number $j$ by $n-j$. If there are $p$ top vertices, the labels are $0,\dots,p-1$ and $n-(n-p),\dots,n$, that is $p,\dots,n$, all distinct. Let $e_m$ join the top vertex $i_m$ and the bottom vertex $j_m$. We claim that the endpoints of $e_m$ are the newest top vertex and the newest bottom vertex among those of $e_1,\dots,e_m$, and that $i_m+j_m=m-1$. For $m=1$ both endpoints are new and $i_1=j_1=0$. For $m\ge2$, the shared endpoint of $e_m$ is an endpoint of $e_{m-1}$, hence by induction the newest vertex of its colour, and it stays so; the other endpoint is new, so it is the newest of its colour and its number is one more than that of the endpoint of $e_{m-1}$ of the same colour. Hence one of $i_m,j_m$ equals its value for $e_{m-1}$ and the other is one larger, and $i_m+j_m=i_{m-1}+j_{m-1}+1=m-1$. The difference along $e_m$ is $(n-j_m)-i_m=n-m+1$, so the differences are $n,n-1,\dots,1$.
\end{proof}

\begin{proposition}[Rosa, 1967]\label{prop:caterpillar}
Every caterpillar is graceful.
\end{proposition}

\begin{proof}
Let the spine be $u_1,\dots,u_s$, and if $s\ge2$ choose $u_1$ and $u_s$ to be the spine ends (a caterpillar with $s=1$ is a star). Colour by parity of distance to $u_1$. Order the edges: the leaf edges at $u_1$, then $u_1u_2$, then the leaf edges at $u_2$, then $u_2u_3$, and so on, ending with the leaf edges at $u_s$. Each leaf edge at $u_i$ shares $u_i$ with the previous edge (which is either the edge $u_{i-1}u_i$ or another leaf edge at $u_i$) and brings a new leaf; each spine edge $u_iu_{i+1}$ shares $u_i$ with the previous edge and brings the new vertex $u_{i+1}$. (If $u_1$ has no leaves, start with $u_1u_2$.) Every vertex appears, since every vertex is on the spine or a leaf of a spine vertex. Lemma~\ref{lem:gracefulorder} applies.
\end{proof}

Not every tree is a caterpillar: the smallest exception is the ``spider'' with three legs of length $2$ (Exercise~\ref{ex:spider}). For such trees no edge order as in Lemma~\ref{lem:gracefulorder} exists, and graceful labellings must be found otherwise.

\begin{remark}[what is known]\label{rem:graceful}
Graceful labellings are known for many families (caterpillars, trees with at most four leaves, all trees with at most $35$ vertices, by a computer search of Fang in 2010, among others); the general conjecture is open. Its decomposition consequence, Ringel's conjecture of 1963 that $K_{2n+1}$ decomposes into copies of any tree with $n$ edges, was proved for all large $n$ by Montgomery, Pokrovskiy and Sudakov (2021), and independently by Keevash and Staden, without graceful labellings.
\end{remark}

\section{Reconstructing a graph from its cards}

Delete one vertex of a graph $G$ on $n\ge3$ vertices, in all $n$ possible ways, and keep the $n$ resulting graphs (the \emph{cards}) up to isomorphism, forgetting which vertex was deleted. The \emph{reconstruction conjecture} of Paul Kelly (1942) and Stanis\l aw Ulam (1960) says that the deck determines $G$ up to isomorphism. Brendan McKay verified it by computer for all graphs on at most $11$ vertices (1997). For a graph $F$, let $s(F,G)$ be the number of subgraphs of $G$ isomorphic to $F$.

\begin{lemma}[Kelly's lemma]\label{lem:kelly}
If $F$ has $f<n$ vertices, then $s(F,G)=\frac1{n-f}\sum_v s(F,G-v)$. In particular $s(F,G)$ is determined by the deck.
\end{lemma}

\begin{proof}
Count pairs (copy of $F$ in $G$, vertex $v$ of $G$ not in the copy). Each copy misses $n-f$ vertices, so there are $(n-f)\,s(F,G)$ pairs. For a fixed $v$, the copies avoiding $v$ are exactly the copies of $F$ in $G-v$, so there are $\sum_v s(F,G-v)$ pairs. The right-hand side depends only on the cards.
\end{proof}

\begin{corollary}\label{cor:kelly}
The number of edges and the degree sequence of $G$ are determined by the deck.
\end{corollary}

\begin{proof}
With $F$ an edge, Kelly's lemma gives the number $m$ of edges. The card $G-v$ has $m-\deg v$ edges, so each card reveals the degree of its deleted vertex.
\end{proof}

From here one can reconstruct regular graphs, disconnected graphs and trees (Kelly), and much more, but not, so far, all graphs.

\section{Lehmer's totient problem}

If $p$ is prime, $\varphi(p)=p-1$ divides $p-1$. In 1932 Derrick Henry Lehmer asked whether a composite $n$ can have $\varphi(n)\mid n-1$. Nobody knows.

\begin{proposition}[Lehmer, 1932]\label{prop:lehmer}
If $n>1$ is composite and $\varphi(n)\mid n-1$, then $n$ is odd, squarefree, has at least three prime factors, and $p-1\mid n-1$ for every prime $p\mid n$.
\end{proposition}

\begin{proof}
If $n>2$, $\varphi(n)$ is even, so $n-1$ is even and $n$ is odd. If $p^2\mid n$, then $p\mid\varphi(n)$ (as $\varphi(p^2)=p(p-1)$ divides $\varphi(n)$ by multiplicativity), so $p\mid n-1$, contradicting $p\mid n$. So $n=p_1\cdots p_r$ with distinct odd primes, and each $p_i-1$ divides $\varphi(n)$, hence $n-1$. Finally suppose $r=2$, $n=pq$ with $3\le p<q$. Then $n-1=(p-1)(q-1)+(p-1)+(q-1)$, so $\varphi(n)=(p-1)(q-1)$ divides $p+q-2>0$. But $(p-1)(q-1)-(p+q-2)=(p-2)(q-2)-1>0$ since $q\ge5$; a contradiction.
\end{proof}

So a counterexample would be a Carmichael number (the last condition is Korselt's criterion) of a very special kind. Lehmer showed it would need at least seven prime factors; Cohen and Hagis (1980) raised this to fourteen. The search for Lehmer numbers is a pleasant exercise in sieving.

\section{Brocard's problem}

In 1876 Henri Brocard asked for which $n$ the number $n!+1$ is a perfect square; Ramanujan asked the same question in 1913. The known solutions are $4!+1=5^2$, $5!+1=11^2$ and $7!+1=71^2$. Berndt and Galway (2000) checked that there are no others with $n\le10^9$, and Overholt (1993) showed that the $abc$ conjecture would imply there are only finitely many.

Why is it hard? A square is a very thin set (there are about $\sqrt N$ squares below $N$), and the heuristic of Chapter~\ref{ch:heuristics} predicts a finite expected number of solutions, $\sum_n(n!)^{-1/2}<\infty$. But the factorials have no structure that interacts with squares in an obvious way, so nothing proves it. The searches rest on a congruence that lets one compute $n!$ modulo a prime quickly when $n$ is close to the prime.

\begin{proposition}\label{prop:wilsonflip}
If $p$ is prime and $0\le n\le p-1$, then $n!\,(p-1-n)!\equiv(-1)^{n+1}\pmod p$.
\end{proposition}

\begin{proof}
By Wilson's theorem, $-1\equiv(p-1)!=n!\prod_{j=n+1}^{p-1}j$. Substituting $j=p-i$ for $i=1,\dots,p-1-n$, the product is $\prod_{i=1}^{p-1-n}(p-i)\equiv(-1)^{p-1-n}(p-1-n)!$. So $n!(p-1-n)!\equiv(-1)^{p-n}=(-1)^{n+1}$, since $p$ is odd (for $p=2$ both sides are $1\equiv-1$).
\end{proof}

So if $n!+1=m^2$, then for each prime $p>n$, $n!+1$ is a square modulo $p$; a sieve over many primes eliminates all $n$ in a range except the true solutions, and the proposition gives $n!\bmod p$ for $n$ close to $p$ from a factorial of a small number.

\section{Sendov's conjecture}

If a polynomial has all its zeros in the closed unit disc, then, conjectured Blagovest Sendov in 1958, every zero has a critical point (a zero of the derivative) at distance at most $1$ from it. The conjecture appeared in Hayman's \emph{Research Problems in Function Theory} (1967), at first attributed to Ilieff. The starting point is a theorem of Gauss and Lucas.

\begin{proposition}[Gauss--Lucas]\label{prop:gausslucas}
The zeros of $p'$ lie in the convex hull of the zeros of $p$, for every nonconstant complex polynomial $p$.
\end{proposition}

\begin{proof}
Write $p(z)=c\prod_{j=1}^n(z-z_j)$. Let $w$ be a zero of $p'$ that is not a zero of $p$ (zeros of both are in the hull trivially). Then $0=\frac{p'(w)}{p(w)}=\sum_j\frac1{w-z_j}=\sum_j\frac{\overline{w-z_j}}{|w-z_j|^2}$. Conjugating, $\sum_j\lambda_j(w-z_j)=0$ with $\lambda_j=|w-z_j|^{-2}>0$, so $w=\sum_j\lambda_jz_j/\sum_j\lambda_j$ is a convex combination of the $z_j$.
\end{proof}

\begin{proposition}\label{prop:sendovweak}
If $p$ has degree $n\ge2$ and all zeros in the closed unit disc, then every zero $a$ of $p$ has a critical point within distance $2n^{-1/(n-1)}$. For $n=2$ this is Sendov's bound $1$.
\end{proposition}

\begin{proof}
Write $p(z)=(z-a)q(z)$ and $p'(z)=nc\prod_{k=1}^{n-1}(z-w_k)$, with $c$ the leading coefficient. Then $p'(a)=q(a)$, so $n|c|\prod_k|a-w_k|=|q(a)|=|c|\prod_{j}|a-z_j|$, the product over the other $n-1$ zeros, each at distance at most $2$ from $a$. Hence $\prod_k|a-w_k|\le\frac{2^{n-1}}n$, and the smallest of the $n-1$ factors is at most $2n^{-1/(n-1)}$. For $n=2$ this is $1$.
\end{proof}

Sendov's conjecture is proved for degree at most $8$ (Brown and Xiang, 1999), and in 2020 Terence Tao proved it for all sufficiently large degrees, with a compactness argument that gives no explicit threshold. The cases in between are open.

\section{How to choose, and where to look}

Some advice, distilled from the stories of this tome and from the people in them.
\begin{itemize}
\item Choose a problem where you can compute. For each problem in this chapter, numerical data can suggest patterns that the statement alone does not reveal.
\item Choose a problem with intermediate questions. ``Prove Singmaster's conjecture'' is a bad research goal; ``find all $a\le10^{30}$ with $N(a)\ge6$'' or ``improve Proposition~\ref{prop:sendovweak} for $n=3$'' are good ones.
\item Read the last three papers on the problem before working on it, and the first one. The first one tells you why people care, and the last ones tell you what is known.
\item Thomas Bloom's database of problems of Paul Erd\H os (\url{https://www.erdosproblems.com}, met in Chapter~\ref{ch:research}) lists hundreds of problems with their status, references and comments, and many have been solved since it opened. It is the best place to look for a problem with a clean statement.
\end{itemize}

\begin{think}
Which of the seven problems would you bet on being solved first? Which partial result of this chapter looks most improvable, and what would the improvement need?
\end{think}

\begin{beauty}
Open problems are the parts of mathematics where the map still says ``here be dragons''. The seven in this chapter are small dragons, which is to say that they are very old and nobody has caught them,
but you can walk right up to them with nothing more than this series and a laptop. That is a privilege few sciences offer.
\end{beauty}

\section{Exercises}

\begin{exercise}[$\star$]\label{ex:singmaster3003}
Find the eight positions of $3003$ in Pascal's triangle, and check Proposition~\ref{prop:singmasterfamily} for $i=2$.
\end{exercise}

\begin{exercise}[$\star\star$]\label{ex:lonelysharp}
Show that for speeds $1,2,\dots,k$, there is no $t$ with $\|it\|>\frac1{k+1}$ for all $i\le k$, so the conjecture is sharp. Find a $t$ achieving $\frac1{k+1}$.
\end{exercise}

\begin{exercise}[$\star$]\label{ex:ringel}
Let $T$ be a gracefully labelled tree with $n$ edges, labels in $\{0,\dots,n\}\subseteq\Z/(2n+1)\Z$. Show that the $2n+1$ translates $T+s$ ($s\in\Z/(2n+1)\Z$), placed on the vertices $\Z/(2n+1)\Z$ of $K_{2n+1}$, use every edge exactly once.
\end{exercise}

\begin{exercise}[$\star\star$]\label{ex:spider}
Let $S$ be the tree with a centre $c$ and three legs $c-a_1-a_2$, $c-b_1-b_2$, $c-d_1-d_2$. Show that $S$ is not a caterpillar, and find a graceful labelling of it.
\end{exercise}

\begin{exercise}[$\star\star$]\label{ex:kellytri}
Show that the number of triangles of $G$ is determined by its deck, and that the deck determines whether $G$ is connected. (For the second part, count the connected cards.)
\end{exercise}

\begin{exercise}[$\star$]\label{ex:sendovdeg2}
Compute, for every zero $a$ of $z^n-1$, the distance from $a$ to the nearest critical point. Deduce that the constant $1$ in Sendov's conjecture cannot be lowered.
\end{exercise}

\section{Hints}

\begin{hint}{ex:singmaster3003}
$3003=\binom{3003}1=\binom{78}2=\binom{15}5=\binom{14}6$ and their mirror images.
\end{hint}

\begin{hint}{ex:lonelysharp}
Place the $k+1$ points $0,t,2t,\dots,kt$ on the circle $\R/\Z$: two of them are within $\frac1{k+1}$ of each other.
\end{hint}

\begin{hint}{ex:ringel}
An edge $\{x,y\}$ of $K_{2n+1}$ has a ``length'' $\min(|x-y|,2n+1-|x-y|)$ in $\{1,\dots,n\}$.
\end{hint}

\begin{hint}{ex:spider}
Removing the leaves leaves a star with three edges. For the labelling, try the centre at $0$.
\end{hint}

\begin{hint}{ex:kellytri}
Triangles: $F=K_3$. Connectedness: $G$ is connected if and only if at least two cards are connected.
\end{hint}

\begin{hint}{ex:sendovdeg2}
The derivative is $nz^{n-1}$.
\end{hint}

\section{Solutions}

\begin{solution}{ex:singmaster3003}
\idea Use the hint and symmetry.
\writeup $3003$ is at $(3003,1)$, $(3003,3002)$, $(78,2)$, $(78,76)$, $(15,5)$, $(15,10)$, $(14,6)$, $(14,8)$: eight positions. For $i=2$: $F_4=3$, $F_6=8$, $F_7=13$, so $n=8\cdot13-1=103$, $k=3\cdot13-1=38$, and $\binom{104}{39}=\binom{103}{40}$, a number with $29$ digits.
\end{solution}

\begin{solution}{ex:lonelysharp}
\idea Pigeonhole on the circle.
\writeup Take any $t$ and the $k+1$ points $0,t,2t,\dots,kt$ modulo $1$. If two coincide, some $\|jt\|=0$ with $1\le j\le k$. Otherwise they cut the circle into $k+1$ arcs, one of length at most $\frac1{k+1}$, between $it$ and $jt$ with $i<j$; then $\|(j-i)t\|\le\frac1{k+1}$ with $1\le j-i\le k$. So $\min_{i\le k}\|it\|\le\frac1{k+1}$ always. The value $t=\frac1{k+1}$ gives $\|it\|=\min(\frac i{k+1},1-\frac i{k+1})\ge\frac1{k+1}$ for $1\le i\le k$.
\end{solution}

\begin{solution}{ex:ringel}
\idea Translation preserves lengths, and a graceful tree has each length once.
\writeup The edge $\{\ell(u),\ell(v)\}$ of the labelled tree has difference $d=|\ell(u)-\ell(v)|\in\{1,\dots,n\}$, which is also its length in $\Z/(2n+1)\Z$, since $d\le n$. The edges of $K_{2n+1}$ of length $d$ are the $2n+1$ pairs $\{x,x+d\}$, and they are exactly the translates of one of them. The tree has exactly one edge of each length $d$, so its $2n+1$ translates contain each edge of length $d$ exactly once, for each $d$. As every edge of $K_{2n+1}$ has some length in $\{1,\dots,n\}$, every edge is used exactly once.
\end{solution}

\begin{solution}{ex:spider}
\idea Remove the leaves and look; then search with the centre fixed.
\writeup The leaves are $a_2,b_2,d_2$; removing them leaves the star with centre $c$ and the three vertices $a_1,b_1,d_1$, which is not a path. So $S$ is not a caterpillar. A graceful labelling (found by a short search): $c=0$, $a_1=1$, $a_2=4$, $b_1=5$, $b_2=3$, $d_1=6$, $d_2=2$. The edge differences are $1,3,5,2,6,4$, that is, $1,\dots,6$ once each.
\end{solution}

\begin{solution}{ex:kellytri}
\idea Kelly's lemma for subgraph counts; a short argument for connectivity.
\writeup Triangles: Kelly's lemma with $F=K_3$ and $n\ge4$. Connectivity: if $G$ is connected, it has a spanning tree, which has at least two leaves, and deleting a leaf of the spanning tree leaves a connected graph; so at least two cards are connected. If $G$ is disconnected with $n\ge3$, a card $G-v$ is connected only if $v$ is an isolated vertex and the rest is connected, and then every other card $G-w$ contains the isolated vertex $v$ and at least one other vertex, so it is disconnected: at most one card is connected. So $G$ is connected if and only if at least two cards are.
\end{solution}

\begin{solution}{ex:sendovdeg2}
\idea Test the bound on the most symmetric examples.
\writeup The zeros of $z^n-1$ are the $n$-th roots of unity, all on the unit circle, and the only critical point is $0$ (with multiplicity $n-1$), at distance exactly $1$ from every zero. So the bound $1$ is attained for every degree $n\ge2$, and no smaller constant can work.
\end{solution}

\section{Sources}

\begin{itemize}
\item D.~Singmaster, How often does an integer occur as a binomial coefficient?, \emph{Amer.\ Math.\ Monthly} 78 (1971), 385--386; Repeated binomial coefficients and Fibonacci numbers, \emph{Fibonacci Quart.} 13 (1975), 295--298; D.~M.~Kane, Improved bounds on the number of ways of expressing $t$ as a binomial coefficient, \emph{Integers} 7 (2007).
\item W.~Bienia, L.~Goddyn, P.~Gvozdjak, A.~Seb\H o, M.~Tarsi, Flows, view obstructions, and the lonely runner, \emph{J.\ Combin.\ Theory Ser.\ B} 72 (1998), with the history from Wills (1967); T.~W.~Cusick, View-obstruction problems, \emph{Aequationes Math.} 9 (1973); J.~Barajas, O.~Serra, The lonely runner with seven runners, \emph{Electron.\ J.\ Combin.} 15 (2008).
\item A.~Rosa, On certain valuations of the vertices of a graph, in \emph{Theory of Graphs} (Rome, 1966), Gordon and Breach, 1967; J.~A.~Gallian, A dynamic survey of graph labeling, \emph{Electron.\ J.\ Combin.}, Dynamic Survey DS6 (updated regularly); R.~Montgomery, A.~Pokrovskiy, B.~Sudakov, A proof of Ringel's conjecture, \emph{Geom.\ Funct.\ Anal.} 31 (2021).
\item P.~J.~Kelly, A congruence theorem for trees, \emph{Pacific J.\ Math.} 7 (1957); S.~M.~Ulam, \emph{A Collection of Mathematical Problems}, Interscience, 1960; B.~D.~McKay, Small graphs are reconstructible, \emph{Australas.\ J.\ Combin.} 15 (1997).
\item D.~H.~Lehmer, On Euler's totient function, \emph{Bull.\ AMS} 38 (1932), 745--751; G.~L.~Cohen, P.~Hagis, On the number of prime factors of $n$ if $\varphi(n)\mid n-1$, \emph{Nieuw Arch.\ Wisk.} 28 (1980).
\item B.~C.~Berndt, W.~F.~Galway, On the Brocard--Ramanujan Diophantine equation $n!+1=m^2$, \emph{Ramanujan J.} 4 (2000); M.~Overholt, The Diophantine equation $n!+1=m^2$, \emph{Bull.\ London Math.\ Soc.} 25 (1993).
\item W.~K.~Hayman, \emph{Research Problems in Function Theory}, Athlone Press, 1967; J.~E.~Brown, G.~Xiang, Proof of the Sendov conjecture for polynomials of degree at most eight, \emph{J.\ Math.\ Anal.\ Appl.} 232 (1999); T.~Tao, Sendov's conjecture for sufficiently high degree polynomials, \emph{Acta Math.} 229 (2022), arXiv:2012.04125.
\end{itemize}
